%!TEX root = ../dissertation.tex

\chapter{Introduzione}
\label{chp:introduzione}

Il calcolo del numero di \textbf{elementi distinti} in un \textbf{flusso di dati} è un problema classico nell'analisi di
grandi moli di dati e nei sistemi che ricevono aggiornamenti in modo continuo.
Applicazioni come monitoraggio di rete, telemetria, osservabilità, analisi di
log, ottimizzazione di interrogazioni e sistemi distribuiti richiedono spesso
di \textit{stimare quanti identificatori diversi siano comparsi in un flusso} \textbf{senza
poter conservare l'intero input in memoria}
\cite{Leskovec_Rajaraman_Ullman_2014,Muthukrishnan_2005,Cormode_2017}. Quando
il \textit{volume dei dati cresce}, il conteggio esatto \textit{diventa costoso} sia in memoria
sia in tempo, mentre il valore cercato deve spesso essere aggiornato con
latenza molto bassa.

Da questa esigenza nasce l'uso di \textbf{strutture riassuntive probabilistiche},
indicate nel seguito come \textbf{sketch}, che comprimono il flusso in uno stato
compatto e consentono di ottenere una stima approssimata del numero di
distinti. La letteratura sul problema si è sviluppata lungo una linea piuttosto
chiara, che parte dal conteggio probabilistico di Flajolet e Martin e arriva a
\textit{HyperLogLog++}, passando per \textit{PCSA}, \textit{LogLog} e \textit{HyperLogLog}
\cite{Flajolet_Martin_1985,Durand_Flajolet_2003,Flajolet_Fusy_Gandouet_Meunier_2007,Heule_Nunkesser_Hall_2013}.
Questa evoluzione migliora progressivamente il compromesso tra occupazione di
memoria, accuratezza e costo di aggiornamento, ma lascia aperto un aspetto che
nei sistemi distribuiti è centrale quanto la stima stessa: \textbf{l'unione di pi\`u sketch}.

Quando un flusso è \textit{suddiviso tra più nodi} o tra più partizioni temporali, la
stima globale non può essere ottenuta ricostruendo ogni volta l'insieme esatto
degli elementi distinti. Diventa quindi \textit{necessario combinare sketch costruiti localmente}.
Nel caso ideale, cioè quando i due lati condividono la stessa struttura, la
stessa semantica di hash e gli stessi parametri, l'operaizone di \textbf{merge} è una proprietà
naturale della struttura riassuntiva. Nei casi reali, però, possono comparire
\textit{eterogeneità} dovute a funzioni di hash diverse, seed diversi oppure precisioni
diverse. In queste situazioni il merge non è più un'operazione banale: occorre
capire quando sia semanticamente corretto, quando sia recuperabile tramite una
normalizzazione esplicita e quando invece produca risultati privi di significato.

Il lavoro sviluppato si concentra proprio su questo punto. Da un lato \textit{vengono
ricostruite e implementate} in modo uniforme alcune delle principali famiglie di
algoritmi per il conteggio approssimato degli elementi distinti. Dall'altro viene
\textit{realizzata un'infrastruttura sperimentale} che separa \textbf{algoritmi}, \textbf{funzioni di
hash}, dataset e raccolta delle metriche, in modo da poter \textit{confrontare i metodi}
in un contesto simile ad un flusso e \textit{studiare l'operazione di merge in scenari distribuiti controllati}.
La parte sperimentale è organizzata in modo da osservare sia il comportamento
della stima lungo il flusso sia quello dell'unione di sketch in caso omogeneo e
nei principali casi eterogenei, con particolare attenzione a \textit{HyperLogLog++} come
riferimento pratico per lo studio dell'operazione di unione.

L'obiettivo non è introdurre un nuovo stimatore teorico, ma \textit{comprendere in modo
sistematico} il comportamento reale degli algoritmi implementati e delle
condizioni di merge. Per fare questo vengono considerati sia il bias sia la
robustezza della stima, viene costruito un dataset sintetico con proprietà
controllate, e vengono osservati separatamente il ruolo della funzione di hash,
del parametro di precisione e della compatibilità semantica tra sketch. In
questo senso il contributo del lavoro è duplice: da un lato conferma in modo
sperimentale proprietà attese della letteratura, dall'altro mette a
disposizione un framework software riutilizzabile per lo studio di nuovi scenari di
merge.

\section{Struttura della tesi}

Il Capitolo~\ref{chp:background} introduce il modello dei flussi di dati, le
funzioni di hash, le strutture riassuntive e le metriche statistiche usate nel
seguito. Il Capitolo~\ref{chp:stato-arte} ricostruisce lo sviluppo storico
degli algoritmi per il conteggio approssimato dei distinti e richiama anche
altre strutture riassuntive impiegate per problemi diversi. Il
Capitolo~\ref{chp:implementazione_nuovo} descrive l'architettura del sistema
sperimentale, le scelte progettuali, gli algoritmi implementati, le funzioni di
hash e il protocollo di generazione dei dati. Il
Capitolo~\ref{chp:risultati_nuovo} presenta i risultati sperimentali, prima in
modalità di flusso e poi negli esperimenti sul merge distribuito in caso
omogeneo ed eterogeneo. Il Capitolo~\ref{chp:conclusioni} raccoglie infine le
considerazioni conclusive, i limiti del lavoro e le principali direzioni di
sviluppo futuro.
